\def\shorttitle{Björner：拓扑方法，第 6、7、9、11 节}
\documentclass[UTF8,zihao=-4,fontset=windows,a4paper]{ctexart}
\usepackage[margin=25mm,headheight=16pt]{geometry}
\usepackage{amsmath,amssymb,amsthm,mathtools,enumitem,fancyhdr,booktabs,longtable,array,xcolor}
\usepackage[unicode,colorlinks,linkcolor=blue!45!black,urlcolor=blue!45!black,bookmarksnumbered]{hyperref}
\setlength{\parskip}{.35em}
\setlength{\emergencystretch}{2em}
\clubpenalty=10000
\widowpenalty=10000
\displaywidowpenalty=10000
\pdfstringdefDisableCommands{\def\quad{ }}
\linespread{1.15}
\setlist{nosep,leftmargin=2em,topsep=.4em}
\pagestyle{fancy}\fancyhf{}\fancyhead[L]{\small 拟阵与单纯复形选读}\fancyhead[R]{\small\shorttitle}\fancyfoot[C]{\thepage}
\newcommand{\sourcepage}[1]{\par\medskip\noindent{\color{gray}\small〔原文第 #1 页〕}\par\smallskip}
\newcommand{\entry}[2]{\par\medskip\noindent\textbf{#1\quad #2}\quad}
\newenvironment{proofcn}{\par\noindent\textbf{证明。}}{\hfill$\square$\par}
\newcommand{\tnote}[1]{\par\smallskip\noindent{\small\color{blue!45!black}\textbf{译注。}#1}\par\smallskip}
\newcommand{\omitsection}[2]{\section*{#1\quad #2}\addcontentsline{toc}{section}{#1\quad #2（留空）}\noindent 本选读稿不译本节；保留节号与标题作为位置标记。\par\vspace{7mm}}
\newcommand{\B}{\mathcal B}\newcommand{\C}{\mathcal C}
\newcommand{\IP}{\operatorname{IP}}\newcommand{\IA}{\operatorname{IA}}\newcommand{\EP}{\operatorname{EP}}\newcommand{\EA}{\operatorname{EA}}
\newcommand{\rk}{\operatorname{rk}}\newcommand{\lk}{\operatorname{lk}}\newcommand{\del}{\operatorname{del}}\newcommand{\con}{\operatorname{con}}
\newcommand{\lex}{\mathrm{lex}}\newcommand{\Gale}{\mathrm{Gale}}
\newcommand{\Q}{\mathbb Q}\newcommand{\R}{\mathbb R}\newcommand{\Z}{\mathbb Z}
\newcommand{\QI}{\mathrm{QI}}\newcommand{\QE}{\mathrm{QE}}\newcommand{\QC}{\mathrm{QC}}
\newcommand{\cm}{\mathrm{CM}}
\newcommand{\real}[1]{\lVert#1\rVert}
\newcommand{\blankrefs}{\section*{参考文献（留空）}\addcontentsline{toc}{section}{参考文献（留空）}正文保留原文的文献引注；完整书目信息请查本目录中的原文 PDF。本选读稿不翻译文末参考文献表。}
\hypersetup{pdfauthor={中文选读译稿}}
\begin{document}

\begin{titlepage}\centering
\vspace*{22mm}{\Huge 拓扑方法\par}\vspace{6mm}{\Large Topological Methods\par}
\vspace{10mm}{\Large Anders Björner\par}\vspace{4mm}瑞典斯德哥尔摩皇家理工学院数学系\par
\vspace{12mm}{\Large 第 6、7、9、11 节中文选译\par}
\vspace{10mm}\begin{minipage}{.87\textwidth}
原文：\emph{Handbook of Combinatorics}，第 34 章（1995），R. Graham、M. Grötschel、L. Lovász 编。以用户提供的 53 页扫描 PDF 为底本。

\textbf{翻译范围：}第 6 节“拟阵与贪心结构”（原书第 1830--1834 页）、第 7 节“有向拟阵”（第 1834--1839 页）、第 9 节“组合拓扑”（原书第 1842--1848 页）和第 11 节“具有特殊结构的复形”（第 1853--1858 页），包括全部编号条目、定理、公式及证明。附译第二部分开头的说明。其余节仅保留标题与位置标记。

\textbf{建议读法：}第 9 节用于查阅定义；第 11 节是本次选读的重点，尤其是 11.2 的可壳性、11.5 的 Cohen--Macaulay 性、11.10(iii) 的拟阵独立复形，以及 11.11--11.14 的诱导子复形结果。被引用但未译的第 10 节可对照原文查阅。

\textbf{记号：}作者采用拓扑学中“不把空集算作面”的单纯复形约定，因此经常写 $\Delta\cup\{\varnothing\}$；$\Delta^k$ 表示 $k$ 维面的集合，$\Delta^{\le k}$ 表示 $k$ 骨架。Samper 一文把空集算作面，并按元素个数记秩，二者相差 1。

\textbf{阅读说明：}本稿保留 1995 年原文对当时研究状态的叙述，不将其视为今日的进展报告。正文的文献引注按原文保留，完整书目请查原文末尾。蓝色“译注”用于补充记号说明或指出原文陈述的适用边界。
\end{minipage}\vfill 2026 年 10 月\par
\end{titlepage}
\tableofcontents\newpage
\omitsection{1}{引言（Introduction）}
\section*{第一部分：例子}\addcontentsline{toc}{section}{第一部分：例子（留空）}
\omitsection{2}{回避性图性质（Evasive graph properties）}
\omitsection{3}{偏序集中的不动点（Fixed points in posets）}
\omitsection{4}{Kneser 猜想（Kneser's Conjecture）}
\omitsection{5}{Borsuk 定理的离散应用（Discrete applications of Borsuk's Theorem）}
\setcounter{section}{5}
\section{拟阵与贪心结构}
\sourcepage{1830}
本节和下一节讨论拟阵及两类相关结构——有向拟阵和贪心结构——的某些拓扑方面。基本定义见 Welsh 的第9章。关于拟阵复形及几何格的其他拓扑事实见(11.10)，另见 Björner (1992)。
\subsection*{基复形与图的划分}
以下结果由 E. Győry 和 L. Lovász 回答 A. Frank 和 S. Maurer 的猜想时证明。
\entry{定理6.1}{（Lovász 1977，Győry 1978）}设 $G=(V,E)$ 是 $k$ 连通图，$\{v_1,\ldots,v_k\}$ 是 $k$ 个顶点，$n_1,\ldots,n_k$ 是正整数且 $\sum_i n_i=|V|$。则存在 $V$ 的划分 $\{V_1,\ldots,V_k\}$，使对每个 $i=1,\ldots,k$，$v_i\in V_i$、$|V_i|=n_i$，且 $V_i$ 在 $G$ 中诱导连通子图。
Lovász 的证明使用拓扑方法，Győry 的证明则没有。本节末将概要介绍 $k=3$ 的 Lovász 证明，以说明拓扑推理的用法。证明依赖于与 $G$ 中某些森林相关的多面体复形的连通性。类似复形可在拟阵的基上、乃至更一般地在贪心结构的基上定义。贪心结构的表述包含前两者，我们用它建立一般结果。先回顾定义。
集合系统 $(E,\mathcal F)$，其中 $\mathcal F\subseteq2^E$，若满足以下公理，称为\textbf{贪心结构}：
\begin{description}
\item[(G1)] $\varnothing\in\mathcal F$；
\item[(G2)] 对每个非空 $A\in\mathcal F$，存在 $x\in A$，使 $A- x\in\mathcal F$；
\item[(G3)] 若 $A,B\in\mathcal F$ 且 $|A|>|B|$，则存在 $x\in A-B$，使 $B\cup x\in\mathcal F$。
\end{description}
若还满足以下条件，则称为\textbf{区间贪心结构}：
\begin{description}\item[(G4)] 若 $A\subseteq B\subseteq C$，其中 $A,B,C\in\mathcal F$，且对某个 $x\in E-C$ 有 $A\cup x,C\cup x\in\mathcal F$，则 $B\cup x\in\mathcal F$。\end{description}
$\mathcal F$ 中的集合称为可行集，极大可行集称为基。所有基的基数相同，记为 $r$，称为贪心结构的秩。
本文仅涉及拟阵（可行集为独立集）和有根图的分枝贪心结构（可行集为构成包含根的树的边集）。二者都是区间贪心结构。其他例子及更多内容见 Welsh 的第9章以及 Korte、Lovász 和 Schrader (1991)、Björner 和 Ziegler (1992) 的综述。
\sourcepage{1831}
除拟阵情形外，贪心结构的可行集不构成单纯复形。不过，按包含关系排序的偏序集 $\overline{\mathcal F}=\mathcal F-\{\varnothing\}$ 的序复形给出了有用的拓扑。称秩为 $r$ 的贪心结构 $(E,\mathcal F)$ 为 $k$ 连通的，如果对每个 $A\in\mathcal F$，存在 $B\in\mathcal F$，使 $A\subseteq B$、$|B-A|=\min(k,r-|A|)$，且每个满足 $A\subseteq C\subseteq B$ 的 $C$ 都属于 $\mathcal F$。拟阵是 $r$ 连通的；$k$ 连通有根图的分枝贪心结构也是 $k$ 连通的。
\entry{命题6.2}{（Björner、Korte 和 Lovász 1985）}设 $(E,\mathcal F)$ 是 $k$ 连通区间贪心结构（$k\ge2$），则其可行集偏序集 $(\overline{\mathcal F},\subseteq)$ 是拓扑意义下的 $(k-2)$ 连通的。
此结果通过定理10.8由(11.10)(iii)推出：$\overline{\mathcal F}$ 的极小元素组成的截集 $C$ 所对应的截集复形 $\Gamma(\overline{\mathcal F},C)$ 是秩至少为 $k$ 的拟阵复形。
设 $\mathcal B$ 为秩 $r$ 的贪心结构的所有基。若 $B_1\cap B_2\in\mathcal F$ 且 $|B_1\cap B_2|=r-1$，称两基相邻。连接所有相邻基，得到以 $\mathcal B$ 为顶点集的\textbf{基图}。基图中最短的圈可明确描述，有两类三角形和一类四边形：
\entry{6.3}{}三个基 $A\cup x,A\cup y,A\cup z$，其中 $A\in\mathcal F$、$|A|=r-1$，张成第一类三角形。
\entry{6.4}{}三个基 $A\cup x\cup y,A\cup x\cup z,A\cup y\cup z$，其中 $A\in\mathcal F$、$|A|=r-2$，张成第二类三角形。
\entry{6.5}{}四个基 $A\cup x\cup u,A\cup x\cup v,A\cup y\cup u,A\cup y\cup v$，其中 $A\in\mathcal F$、$|A|=r-2$，张成一个四边形。
分枝贪心结构不会出现第二类三角形。现在在每个三角形及四边形内贴上一个二维胞腔（“膜”），得到二维正则胞腔复形 $\mathcal K$，称为\textbf{基复形}。不难检验，秩不超过2的任何2连通贪心结构，其基复形都是1连通的，即连通且单连通。秩为2时，这直接由交换公理(G3)得到。较高秩时有：
\entry{定理6.6}{（Björner、Korte 和 Lovász 1985）}任何3连通区间贪心结构的基复形 $\mathcal K$ 都是1连通的。
为了说明第二部分的一些工具，给出简短证明。设 $P$ 是 $\mathcal K$ 的闭胞腔按包含关系组成的偏序集，$Q$ 是 $(\mathcal F,\subseteq)$ 的最高三层，即秩分别为 $r-2,r-1,r$ 的可行集。令 $f:P\to Q$ 为反序映射，把每个胞腔映到张成它的各个基的交。由命题6.2和引理11.12，$Q$ 是1连通的；由定理10.5，只需检验所有纤维 $f^{-1}(Q_{\ge A})$ 是1连通的。
\sourcepage{1832}
若 $r(A)=r-i$，$i=0,1,2$，则 $f^{-1}(Q_{\ge A})$ 是收缩 $A$ 所得秩为 $i$ 的贪心结构的基复形；我们已检验秩不超过2的情形是1连通的。
设 $P=B_1B_2\cdots B_d$、$Q=B_dB_{d+1}\cdots B_g$ 是拟阵基图中的路，令 $PQ=B_1\cdots B_g$ 为其拼接。若 $R$ 形如 $BCB$、$BCDB$ 或 $BCDEB$，其中 $B=B_d$，则称 $PQ$ 与 $PRQ$ 相差一个\textbf{初等同伦}。
\entry{定理6.7}{（Maurer 1973）}拟阵基图中具有相同端点的任意两条路 $P,P'$，可通过一系列初等同伦相互变换。
Maurer 的同伦定理显然是定理6.6在拟阵情形下的组合表述。下一节将给出其在有向拟阵中的应用。
现在回到定理6.1。下面引用 Lovász (1979) 对 $k=3$ 的证明概要，方括号中作了若干调整以适应本文记号。
设 $G$ 是3连通图，$v_1,v_2,v_3\in V(G)$，且 $n_1+n_2+n_3=|V(G)|$。加入新点 $a$ 并将其与 $v_1,v_2,v_3$ 相连。考虑为新图 $G'$ 构造的拓扑空间 $\mathcal K$。［按本文语言，$\mathcal K$ 是有根图 $(G',a)$ 所确定的分枝贪心结构的基复形；这个贪心结构是3连通的，其基是 $G'$ 的生成树。］对每棵生成树 $T$，令 $f_i(T)$ 为从 $a$ 沿边 $(a,v_i)$ 在 $T$ 中可达的点数（$i=1,2$）。映射 $f:T\mapsto(f_1(T),f_2(T))$ 将 $\mathcal K$ 的顶点映为平面格点。对每个四边形二维胞腔加一条对角线，将其分成两个三角形，得到三角剖分 $\widetilde{\mathcal K}$。在每个三角形上线性延拓 $f$，得到到平面的连续映射。其像显然包含于三角形 $\Delta=\{x\ge0,y\ge0,x+y\le n\}$；我们证明该映射覆盖 $\Delta$。
先取生成树 $T_1,T_2,T_3$，使 $f(T_1)=(n,0)$、$f(T_2)=(0,n)$、$f(T_3)=(0,0)$。应用2连通贪心结构的基图连通这一事实，在 $G'-(a,v_3)$ 中选取连接 $T_1,T_2$ 的折线 $P_{12}$，使其所有点满足 $f_3=0$。于是 $f(P_{12})$ 沿 $\Delta$ 的边连接 $(n,0)$ 与 $(0,n)$。类似定义 $P_{23},P_{31}$。
由定理6.6，$P_{12}+P_{23}+P_{31}$ 可在 $\mathcal K$ 中收缩为一点。因此其像也可在 $f(\mathcal K)$ 中收缩为一点。但三角形边界不能在去掉一个内部点的三角形中收缩为一点（Brouwer 定理13.1），故 $f(\mathcal K)$ 必覆盖整个 $\Delta$。特别地，$(n_1,n_2)$ 属于某个三角形的像。容易看出，这意味着它属于某个顶点的像：即存在生成树 $T$ 满足 $f_1(T)=n_1$、$f_2(T)=n_2$。
\sourcepage{1833}
于是 $T-a$ 的三个分支给出所需的 $V(G)$ 的划分。
定理6.6是更一般结果的特例：对任意 $k$ 连通区间贪心结构，某个高维基复形是 $(k-2)$ 连通的。将上述 $k=3$ 的思想推广，即得一般的定理6.1。详情见 Lovász (1977) 和 Björner、Korte、Lovász (1985)。
\subsection*{Tutte 同伦定理}
若拟阵能在每个域上坐标化，称其为\textbf{正则拟阵}。Tutte (1958) 用排除子式刻画了正则拟阵，其证明本质上依赖一个“同伦定理”：某些二维复形是1连通的。R. Reid 和 R. Bixby 也使用它证明了在 $GF(3)$ 上可表示性的排除子式刻画。后来 P. Seymour 等人找到了不使用此定理的其他证明；见 Seymour 的第10章。
Tutte 同伦定理似乎是组合学中这类拓扑结果最早的例子。完整叙述及证明都相当繁琐，这里只陈述足以理解其性质的内容；详情见 Tutte (1958, 1965)。
设 $L$ 是秩 $r$ 的有限几何格，$L^i$ 是秩为 $i$ 的平坦集之集；$L^{r-1},L^{r-2},L^{r-3}$ 的元素分别称为余点、余线、余平面。把平坦集 $X$ 视为点集 $L^1$ 的子集 $\bar X=\{p\in L^1:p\le X\}$。给定点 $a\in L^1$，在顶点集 $L^{r-1}_{\not\ge a}=\{X\in L^{r-1}:X\not\ge a\}$ 上定义图 $TG(L,a)$：两个不含 $a$ 的余点 $X,Y$ 相邻，当且仅当 $X\wedge Y$ 是余线且 $\bar X\cup\bar Y\ne L^1-a$。在该图的下列三角形和四边形中贴上二维胞腔，构造二维正则胞腔复形 $TC(L,a)$：
\entry{6.8}{}满足 $\rk(X\wedge Y\wedge Z)\ge r-3$ 的三角形 $XYZX$。
\entry{6.9}{}四边形 $XYZTX$，其中 $P=X\wedge Y\wedge Z\wedge T$ 满足 $\rk(P)=r-3$，且余线 $P\vee a$ 恰被两个余点覆盖，或者区间 $[P,\hat1]$ 同构于从 Fano 拟阵 $F_7$ 中删去一个点所得拟阵的平坦集格。
若 $L$ 没有同构于 Fano 拟阵对偶 $F_7^*$ 的子式，则(6.8)、(6.9)描述了 Tutte 复形 $TC(L,a)$ 的所有二维胞腔。［这意味着在表示理论的应用中，上述定义已足够。］一般还需要在某些满足 $\rk(X\wedge Y\wedge Z\wedge T)=r-4$ 的四边形上贴二维胞腔；这种“余秩4型”四边形的定义相当复杂，这里不再描述。
\sourcepage{1834}
\entry{定理6.10}{（同伦定理，Tutte 1958）}复形 $TC(L,a)$ 是1连通的。
其组合含义是：不含 $a$ 的任意两个余点 $X,Y$ 可在 Tutte 图 $TG(L,a)$ 中由不含 $a$ 的路连接；任意两条这样的路相差一系列初等同伦，其形式为 $XYX$、(6.8)中的 $XYZX$、(6.9)中的 $XYZTX$ 或余秩4型。比较定理6.7前的讨论。
此处的表述与 Tutte (1958) 原表述形式不同，内容相同。据 Tutte (1979，第446页)，证明虽长，但本质上是图论和几何性的。除特殊情形外，似乎没有显著简化原证明的方法。其中一种特殊情形是：对所有不含 $a$ 且交为余线的余点 $X,Y$，都有 $\bar X\cup\bar Y\ne L^1-a$。此时由(11.10)(iv)、(11.2)及定理11.14，$L-[a,\hat1]$ 的最高三层构成1连通偏序集；类似定理6.6的证明，应用纤维定理10.5，可将1连通性传给 $TC(L,a)$。寻找一般情形下更简单、更概念化的证明仍有意义。
篇幅不允许详细说明定理6.10在表示理论中的应用，只概要介绍思想。Tutte 正则拟阵刻画的充分性证明对底集大小归纳，故需要删除点 $a$。粗略说，“正则”坐标化定义在余点上，新点 $a$ 处的值沿 Tutte 图中的路从一个余点传到另一个。此时同伦定理保证不同路径不会产生矛盾。定理7.6的证明将更详细展示类似思想。

\section{有向拟阵}
本节讨论有向拟阵理论中的两个主题。最重要的是 Folkman 和 Lawrence (1978) 的\textbf{拓扑表示定理}：每个有向拟阵都能由伪球面配置实现。作为应用，我们说明这类实现如何快速证明秩为3的有向拟阵的某些组合性质。其次，我们依照 Las Vergnas (1978)，概要说明怎样用 Maurer 同伦定理6.7推出行列式符号函数的存在性。
有向拟阵的定义见 Welsh 的第9章。由于这里使用 Folkman 和 Lawrence (1978) 的略有不同的表述，并要用线性情形作为动机，先简要回顾基础并固定记号。更详细的讨论见 Bachem 和 Kern (1992)，以及 Björner、Las Vergnas、Sturmfels、White 和 Ziegler (1993)。
\sourcepage{1835}
设 $E$ 是有限集，带有无不动点的对合 $x\mapsto x^*$，即对所有 $x\in E$，$x^*\ne x=x^{**}$。对 $A\subseteq E$ 写 $A^*=\{x^*:x\in A\}$。\textbf{有向拟阵} $\mathcal O=(E,*,\mathcal C)$ 由这样的集合及一个非空子集族 $\mathcal C$ 构成，满足：
\begin{description}
\item[(OM1)] $\mathcal C$ 是杂族，即不同的 $C_1,C_2\in\mathcal C$ 满足 $C_1\nsubseteq C_2$；
\item[(OM2)] 若 $C\in\mathcal C$，则 $C^*\in\mathcal C$ 且 $C\cap C^*=\varnothing$；
\item[(OM3)] 若 $C_1,C_2\in\mathcal C$、$C_1\ne C_2^*$ 且 $x\in C_1\cap C_2^*$，则存在 $D\in\mathcal C$，使 $D\subseteq C_1\cup C_2-\{x,x^*\}$。
\end{description}
$\mathcal C$ 中的集合称为 $\mathcal O$ 的回路。对 $x\in E$ 令 $\bar x=\{x,x^*\}$；对 $A\subseteq E$ 令 $\bar A=\{\bar x:x\in A\}$，并令 $\overline{\mathcal C}=\{\bar C:C\in\mathcal C\}$。$\overline{\mathcal C}$ 满足通常的拟阵回路交换公理，因此 $\overline{\mathcal O}=(\bar E,\overline{\mathcal C})$ 是拟阵，称为 $\mathcal O$ 的\textbf{底层拟阵}。并非所有拟阵都能这样得到；能这样得到的称为可定向的。若 $\bar B$ 是底层拟阵的基，称 $B\subseteq E$ 是 $\mathcal O$ 的基。$\mathcal O$ 的秩等于底层拟阵的秩。以下不妨默认所有有向拟阵都是简单的，即没有少于三个元素的回路。
有向拟阵的基本模型是 $\R^d$ 中的向量集及正线性相关关系，更一般地也可使用任意有序域上的向量。设有限集 $E\subseteq\R^d-\{0\}$ 满足 $E=-E$，且 $E$ 中不同的平行向量必互为相反数。令 $x^*=-x$。若存在不全为零的实系数 $\lambda_x\ge0$，使 $\sum_{x\in A}\lambda_xx=0$，称 $A\subseteq E$ 正线性相关。令 $\mathcal C$ 为按包含关系极小的正线性相关集，但排除形如 $\{x,x^*\}$ 的集合。等价地，$\mathcal C$ 由这样的集合组成：它们是一个维数至少为2的单形的顶点集，且原点属于该单形的相对内部。以此方式得到的有向拟阵称为实数域上的\textbf{线性有向拟阵}或\textbf{可实现有向拟阵}。并非所有有向拟阵都同构于线性的有向拟阵。
\subsection*{拓扑表示定理}
为了引出表示定理，先考虑线性情形。实际上，表示定理说：从线性情形得到的直观，在允许适当的拓扑变形后，对一般有向拟阵仍然基本正确。
设有限集 $E\subseteq\R^d-\{0\}$ 满足 $E=-E$，$\mathcal O=(E,*,\mathcal C)$ 是前述线性有向拟阵。对每个 $e=\bar x=\{x,x^*\}\in\bar E$，令 $H_e$ 为与 $e$ 所张成直线正交的超平面。超平面配置 $\mathcal H=\{H_e:e\in\bar E\}$ 包含 $\mathcal O$ 的全部信息：由 $H_e$ 可恢复一对方向相反的法向量，而 $\mathcal O$ 的回路定义与向量长度无关。与单位球面 $S^{d-1}$ 相交后，也可考虑球面配置 $\mathcal S=\{H_e\cap S^{d-1}:e\in\bar E\}$，即 $(d-1)$ 维球面中的一族赤道 $(d-2)$ 维球面。因此线性有向拟阵（允许重新定向）、超平面配置及球面配置是等价的。
\sourcepage{1836}
把线性有向拟阵看成球面配置时，宜将元素 $x\in E$ 想成闭半球 $\bar H_x=\{y\in S^{d-1}:\langle y,x\rangle\ge0\}$。于是 $A\in\mathcal C$ 当且仅当 $A\cap A^*=\varnothing$，且 $A$ 是满足 $\bigcup_{x\in A}\bar H_x=S^{d-1}$ 的极小集合。
需要以下拓扑术语。若拓扑空间 $\Sigma$ 与标准 $j$ 维球面 $S^j=\{x\in\R^{j+1}:\|x\|=1\}$ 同胚，则称之为球面。$\Sigma$ 中的\textbf{伪球面} $S$ 是赤道 $\{x\in S^j:x_{j+1}=0\}$ 在某个同胚 $f:S^j\to\Sigma$ 下的像；拓扑学中也称这种嵌入为温驯或平坦的余维一子球面。其两侧（伪半球）为 $S^+=f(\{x\in S^j:x_{j+1}\ge0\})$、$S^-=f(\{x\in S^j:x_{j+1}\le0\})$。显然 $S$ 是两侧之交，两侧均同胚于球。
关键定义是：$S^{d-1}$ 中的\textbf{伪球面配置} $(\bar E,\mathcal A)$ 是一族有限个不同的伪球面 $\mathcal A=\{S_e:e\in\bar E\}$，满足：
\begin{description}
\item[(AP1)] 每个非空交 $S_A=\bigcap_{e\in A}S_e$（$A\subseteq\bar E$）都是球面；
\item[(AP2)] 对每个非空交 $S_A$ 和每个 $e\in\bar E$，或者 $S_A\subseteq S_e$，或者 $S_A\cap S_e$ 是 $S_A$ 中的伪球面，其两侧为 $S_A\cap S_e^+$ 和 $S_A\cap S_e^-$。
\end{description}
这一定义来自 Folkman 和 Lawrence (1978)。他们实际上要求了更多条件，但 Mandel (1982) 证明额外条件是多余的。
仿照线性情形，伪球面配置 $\mathcal A$ 给出系统 $\mathcal O(\mathcal A)=(E,*,\mathcal C)$：令 $E=\{S_e^+:e\in\bar E\}\cup\{S_e^-:e\in\bar E\}$，令 $(S_e^+)^*=S_e^-$，反之亦然，并定义 $\mathcal C$ 为满足 $\bigcup A=S^{d-1}$ 且 $A\cap A^*=\varnothing$ 的极小子集 $A\subseteq E$ 的集合。尽管存在拓扑变形，$\mathcal O(\mathcal A)$ 仍是有向拟阵；更令人惊讶的是，这样能得到所有有向拟阵。若 $\bigcap\mathcal A=\varnothing$，称该配置为\textbf{本质的}。
\entry{定理7.1}{（表示定理，Folkman 和 Lawrence 1978）}\begin{enumerate}[label=(\roman*)]
\item 若 $\mathcal A$ 是 $S^{d-1}$ 中的伪球面配置，则 $\mathcal O(\mathcal A)$ 是有向拟阵；若配置本质，则其秩为 $d$。
\item 若 $\mathcal O$ 是秩 $d$ 的有向拟阵，则存在 $S^{d-1}$ 中的本质伪球面配置 $\mathcal A$，使 $\mathcal O=\mathcal O(\mathcal A)$。
\item 在自然等价关系下，映射 $\mathcal A\mapsto\mathcal O(\mathcal A)$ 给出秩 $d$ 的有向拟阵与 $S^{d-1}$ 中本质伪球面配置之间的一一对应。
\end{enumerate}
这个结果的证明相当复杂。对(ii)，先由有向拟阵构造偏序集，再由定理12.6证明该偏序集是某个正则胞腔复形 $\mathcal C$ 的面偏序集。这个复形给出 $(d-1)$ 维球面及构成配置的各种 $(d-2)$ 维子球面。$\mathcal C$ 是可构造的（Edmonds 和 Mandel 1978；Mandel 1982），甚至是可壳的（Lawrence 1984）；因此整个构造可在分片线性拓扑中进行，不必处理病态拓扑现象。完整证明见 Folkman 和 Lawrence (1978)、Mandel (1982) 及 Björner 等 (1993)。
\sourcepage{1837}
表示定理说明，秩为3的有向拟阵对应于二维球面上的“伪圆”配置；在射影版本中则对应实射影平面上的伪直线配置。这种表示能迅速证明某些组合性质：
\entry{定理7.2}{}设 $M$ 是秩3的可定向拟阵。则：(i) $M$ 有一条含两个点的直线；(ii) 若将 $M$ 的点染成两色，则存在一条单色直线。
下面说明如何由定理7.1得到定理7.2。把 $M$ 的点表示为二维球面上的伪圆；直线对应于交非空的极大伪圆族，其交必为零维球面，即两个点。伪圆配置确定图 $G$：顶点是交点，边是相邻交点间的伪圆段。$G$ 嵌入 $S^2$，因此是平面图；又因 $\rk(M)=3$，它是简单图。需要以下引理。
\entry{引理7.3}{}对任意已作平面嵌入的简单图：(i) 存在度数至多为5的顶点；(ii) 若边染成两色，则存在一个顶点，使每种颜色的边在嵌入所诱导的循环次序中都连续排列。
(i) 是欧拉公式的著名推论；(ii) 也是欧拉公式的推论，但较不为人熟知。Cauchy 在证明三维凸多胞体刚性定理时曾使用它。
现在完成定理7.2的证明。若 $M$ 的每条直线至少有3个点，则伪圆配置所确定的图 $G$ 的每个顶点度数至少为6，与(i)矛盾。若伪圆染成两色，且每个交点都至少有一条每种颜色的伪圆通过，则诱导的边染色将违反(ii)。
定理7.2第一部分是 Sylvester--Gallai 定理的推广，线性情形的证明在20世纪40年代已有。Csima 和 Sawyer (1993) 进一步证明：$M$ 中两点直线至少有 $\frac6{13}|M|$ 条，也使用了伪直线表示。第二部分在线性情形下的证明由 G. D. Chakerian 给出，后由 Edmonds、Lovász 和 Mandel (1980) 重新发现，他们也观察到了对有向拟阵的推广。
\sourcepage{1838}
\subsection*{基的符号赋值}
正如普通拟阵，有向拟阵也可用多种方式刻画。这里讨论有向拟阵的基集 $\mathcal B$ 的一个特征性质：可以定义类似行列式的符号，但不定义其大小。Las Vergnas (1978) 最先证明此事；J. Bokowski、A. Dress、L. Gutierrez-Novoa 和 J. Lawrence 也发现了用带符号的基刻画有向拟阵的方法。
考虑函数 $\delta:\widetilde{\mathcal B}\to\{+1,-1\}$：它将线性有向拟阵 $(E,*,\mathcal C)$（$E\subseteq\R^d$）的有序基映为其行列式的符号。对 $\R^d$ 中某些有序基对 $\beta,\beta'$，可如下定义函数 $\eta$：
\entry{7.4}{}若 $\beta,\beta'$ 是同一个基 $B$ 的排列，则同奇偶性时令 $\eta(\beta,\beta')=+1$，否则令其为 $-1$。
\entry{7.5}{}若 $\beta=x_1\cdots x_{r-1}y$、$\beta'=x_1\cdots x_{r-1}z$ 且 $y\ne z$，当 $y,z$ 位于 $\{x_1,\ldots,x_{r-1}\}$ 张成的超平面的同侧时，令 $\eta(\beta,\beta')=+1$；否则令其为 $-1$。
选定有序基 $\beta_0$ 并令 $\det(\beta_0):=+1$ 后，由线性代数规则可确定所有有序基的行列式及其符号 $\delta$。但 $\delta$ 也由组合信息决定：任意两个有序基可由(7.4)或(7.5)型步骤组成的链连接，而对这样相关的 $\beta,\beta'$，有 $\delta(\beta)=\eta(\beta,\beta')\delta(\beta')$。
上面的讨论指出如何把行列式符号函数推广到所有有向拟阵。首先，为使(7.5)与公理(OM1)--(OM3)更相容，改写如下：
\entry{7.5'}{}设 $\beta=x_1\cdots x_{r-1}y$、$\beta'=x_1\cdots x_{r-1}z$ 且 $y\ne z$。若 $y\ne z^*$，令 $\{C,C^*\}$ 为唯一一对回路，使底层拟阵中 $\{\bar y,\bar z\}\subseteq\bar C\subseteq\{\bar x_1,\ldots,\bar x_{r-1},\bar y,\bar z\}$。若 $y,z$ 中一个属于 $C$、另一个属于 $C^*$，令 $\eta(\beta,\beta')=+1$，否则令其为 $-1$。
\tnote{原文在(7.5')中没有单独写出 $y=z^*$ 时的约定；在线性情形，它们位于超平面两侧，因此取 $\eta=-1$。}
\entry{定理7.6}{（Las Vergnas 1978）}设 $\widetilde{\mathcal B}$ 是有向拟阵的有序基集合，$\beta_0\in\widetilde{\mathcal B}$。存在唯一函数 $\delta:\widetilde{\mathcal B}\to\{+1,-1\}$，使 $\delta(\beta_0)=+1$，且若 $\beta,\beta'$ 按(7.4)或(7.5')相关，则 $\delta(\beta)=\eta(\beta,\beta')\delta(\beta')$。
证明思路如下。在顶点集 $\widetilde{\mathcal B}$ 上定义图：若一对有序基按(7.4)或(7.5')相关，则连边。此图显然连通，且有到其底层拟阵基图 $\mathcal B$ 的投影 $\pi:\widetilde{\mathcal B}\to\mathcal B$。令 $\delta(\beta_0)=+1$，并对 $\beta\in\widetilde{\mathcal B}$，选取从 $\beta_0$ 到 $\beta$ 的路 $\beta_0,\beta_1,\ldots,\beta_n=\beta$，定义
\[\delta(\beta):=\prod_{i=1}^n\eta(\beta_{i-1},\beta_i).\]
\sourcepage{1839}
只需证明定义与所选路径无关。若 $P_1,P_2$ 是两条这样的路，则由定理6.7，它们在基图中的投影 $\pi(P_1),\pi(P_2)$ 相差一系列初等同伦。因此只需对有序基图中投影为边的往返路 $BCB$、三角形 $BCDB$ 或四边形 $BCDEB$ 的闭路 $\alpha_0,\ldots,\alpha_k=\alpha_0$，验证
\[\prod_{i=1}^k\eta(\alpha_{i-1},\alpha_i)=1.\]
而产生基图三角形和四边形的基配置已在(6.3)--(6.5)中明确刻画，故验证被化为可处理的有限情形。详情见 Las Vergnas (1978)。

\omitsection{8}{Hard Lefschetz 定理的离散应用}
\clearpage
\section*{第二部分：工具}\addcontentsline{toc}{section}{第二部分：工具}
\sourcepage{1842}
本章余下部分回顾组合拓扑中的一些定义与结果，它们已经证明在组合学中尤其有用。第 9 节（单纯复形）、第 12 节（胞腔复形）和第 13 节（不动点与对径性定理）的内容十分一般，许多拓扑学书籍都有详细论述，因此只零星给出具体参考文献。相比之下，第 10、11 节的大多数主题更为专门，将给出更充分的文献，有时也给出证明。所提到的许多结果已在大量论文和书籍中讨论过。适当时，我们引用所知的原始出处，以及给出简洁证明、推广或较新讨论的一些论文；这样的选择难免带有主观性。对于可能无意造成的不准确或遗漏，谨致歉意。
\setcounter{section}{8}
\section{组合拓扑}
本节回顾单纯复形的基本事实。Munkres (1984a) 和 Spanier (1966) 是很好的通用参考书。假定读者已了解（拓扑）空间、连续映射、同胚等基本概念。本章中，拓扑空间之间的映射一律假定连续，即使没有明说。
\subsection*{单纯复形与偏序集}\addcontentsline{toc}{subsection}{单纯复形与偏序集}
\entry{9.1}{}一个\textbf{抽象单纯复形} $\Delta=(V,\Delta)$ 由集合 $V$（顶点集）及其非空有限子集组成的一族 $\Delta$ 构成；这些子集称为\textbf{单形}或\textbf{面}，满足
\[
\varnothing\ne\sigma\subseteq\tau\in\Delta\quad\Longrightarrow\quad\sigma\in\Delta.
\]
通常 $V=\bigcup_{\tau\in\Delta}\tau$，于是可以在记号中略去 $V$。面的\textbf{维数}为 $\dim\sigma=|\sigma|-1$，复形的维数为 $\dim\Delta=\max_{\sigma\in\Delta}\dim\sigma$。若每个面都包含于某个 $d$ 维面，则称 $d$ 维复形为\textbf{纯的}。由一个 $d+1$ 元集合的全部非空子集组成的复形，称为 $d$ \textbf{单形}。定义允许空复形 $\Delta=\varnothing$，约定其维数为 $-1$。

\textbf{原文附注：}这里采用不含空面的定义，是拓扑学中的标准定义。在组合学中，通常容许空集作为面更方便，尤其与约化同调的定义一致。

\sourcepage{1843}
记 $\Delta^k=\{\Delta\text{ 的 }k\text{ 维面}\}$，且对 $k\ge0$，令 $\Delta^{\le k}=\bigcup_{j\le k}\Delta^j$。$\Delta^0=V$ 和 $\Delta^1$ 的元素分别称为\textbf{顶点}与\textbf{边}。若 $\Delta$ 纯且为 $d$ 维，$\Delta^d$ 的元素称为\textbf{极大面}（facet，亦称 chamber）。$\Delta^{\le k}$ 是 $\Delta$ 的\textbf{$k$ 骨架}，它是一个子复形。

\textbf{几何单纯复形}是 $\R^d$ 中的多面体复形〔按 12.1 的意义〕，其胞腔为几何单形，即仿射无关点集的凸包。若 $\Gamma$ 是几何单纯复形，其胞腔极点集组成有限抽象单纯复形 $\Delta(\Gamma)$。反过来，对有限 $d$ 维抽象单纯复形 $\Delta$，在 $\R^{2d+1}$ 中存在几何单纯复形 $\Gamma$，使 $\Delta(\Gamma)\cong\Delta$。这种 $\Gamma$ 的底空间 $\bigcup\Gamma$ 在线性同胚意义下唯一，称为 $\Delta$ 的\textbf{几何实现}或\textbf{空间}，记为 $\real\Delta$。反过来，$\Delta$ 称为 $\real\Delta$ 以及任何与其同胚空间的一个\textbf{三角剖分}。因此，在有限情形，更一般地在有限维、可数且局部有限的情形，抽象与几何单纯复形是等价的概念。任意无限抽象单纯复形的几何实现可按 Spanier (1966) 构造。

\textbf{单纯映射} $f:\Delta_1\to\Delta_2$ 是顶点映射 $f:\Delta_1^0\to\Delta_2^0$，使对全部 $\sigma\in\Delta_1$ 都有 $f(\sigma)\in\Delta_2$。在各单形上仿射延拓，便诱导连续映射 $\real f:\real{\Delta_1}\to\real{\Delta_2}$。

所有 $d$ 维单纯复形在 $\R^{2d+1}$ 中存在直线实现很容易证明，而且 $2d+1$ 是最佳的一般界。然而，在特殊情形中能否在 $\R^j$（$d<j\le2d$）内进行直线实现或拓扑实现，是困难且研究很多的问题。$d=1$ 时这是图的平面性问题（见 Thomassen 所写的第 5 章）；$d\ge2$ 的直线嵌入见例如 Bokowski 与 Sturmfels (1989) 及其文献，拓扑嵌入见 Rushing (1973)。例如，原文写作时尚不知道二维环面的每个三角剖分是否都能直线嵌入 $\R^3$。经典的 van Kampen--Flores 定理（1932--1933）断言：一个 $(2d+2)$ 单形的 $d$ 骨架不能嵌入 $\R^{2d}$。Sarkaria (1991b) 在也包含第 5 节所述拓扑 Radon--Tverberg 定理的框架中讨论了这个结果；另见 Sarkaria (1991a)。

\entry{9.2}{}设 $P=(P,\le)$ 为\textbf{偏序集}。全序子集 $x_0<x_1<\cdots<x_k$ 称为长度 $k$ 的\textbf{链}。所有链长的上确界称为 $P$ 的\textbf{秩}或\textbf{长度}。若所有极大链都有相同的有限长度，称 $P$ 为\textbf{纯的}。若每对 $x,y\in P$ 都有最小上界（并）$x\vee y$ 和最大下界（交）$x\wedge y$，则称 $P$ 为\textbf{格}。记
\[
P_{\ge x}=\{y\in P:y\ge x\},
\]
其他类似记号同理。对 $x<y$，定义\textbf{开区间} $(x,y)=P_{>x}\cap P_{<y}$ 及\textbf{闭区间} $[x,y]=P_{\ge x}\cap P_{\le y}$。满足 $\hat0\le x$、$x\le\hat1$ 对全部 $x\in P$ 成立的元素，分别称为\textbf{最小元}和\textbf{最大元}。二者均存在时称 $P$ \textbf{有界}。

\sourcepage{1844}
此时记 $\overline P=P\setminus\{\hat0,\hat1\}$，称为 $P$ 的\textbf{真部}。对任意偏序集 $P$，$\widehat P=P\cup\{\hat0,\hat1\}$ 表示添加新的最大、最小元后的偏序集，故 $|\widehat P\setminus P|=2$。

若 $P$ 纯且秩为 $r$，令 $r(x)=\rk(P_{\le x})$。\textbf{秩函数} $r:P\to\{0,1,\ldots,r\}$ 在每条极大链上为双射。它将 $P$ 分成秩层 $P_i=\{x\in P:r(x)=i\}$，$0\le i\le r$。

\entry{9.3}{}单纯复形 $\Delta$ 的\textbf{面偏序集} $\mathcal P(\Delta)=(\Delta,\subseteq)$ 是按包含关系排序的面集。\textbf{面格}是 $\widehat{\mathcal P}(\Delta)=\mathcal P(\Delta)\cup\{\hat0,\hat1\}$，它确为格。$\mathcal P(\Delta)$ 纯当且仅当 $\Delta$ 纯，而且 $\rk\mathcal P(\Delta)=\dim\Delta$。

偏序集 $P$ 的\textbf{序复形} $\Delta(P)$ 以 $P$ 为顶点集，其 $k$ 维面是 $P$ 的 $k$ 链 $x_0<\cdots<x_k$。保序映射或逆序映射 $f:P_1\to P_2$ 都给出单纯映射 $\Delta(P_1)\to\Delta(P_2)$，因而诱导连续映射 $\real f:\real{\Delta(P_1)}\to\real{\Delta(P_2)}$。$\Delta(P)$ 的定义可追溯到 Aleksandrov (1937)。

复形 $\Delta$ 的 $\operatorname{sd}\Delta:=\Delta(\mathcal P(\Delta))$ 称为（第一次）\textbf{重心细分}，其名称来自几何版本。基本事实是 $\Delta$ 与 $\operatorname{sd}\Delta$ 同胚。因此，经由 $\mathcal P(\cdot)$ 与 $\Delta(\cdot)$ 在单纯复形和偏序集之间转换，不改变拓扑；从拓扑角度看，两种概念本质上等价。

定义 $\real P:=\real{\Delta(P)}$，便为每个偏序集配上一个拓扑空间。本章中关于 $P$ 的拓扑陈述，一律是指这个空间。

至少还有另一种给偏序集配备有用拓扑的方式，也来自 Aleksandrov (1937)：把序理想作为开集。序理想是满足 $x\le y\in A\Rightarrow x\in A$ 的子集 $A\subseteq P$。将所得空间记为 $T(P)$。例如，第 12 节图 2 右边的偏序集，其 $T(P)$ 恰有十个开集，而其序复形同胚于二维球面。对这种理想拓扑，连续映射恰为保序映射，而同伦〔见 9.10〕具有直接的组合意义。例如，$T(P)$ 可缩当且仅当 $P$ 按 11.1 的意义可拆除，见 Stong (1966)。理想拓扑与偏序集上的层上同调有关（Baclawski 1975；Yuzvinsky 1987），并与序复形拓扑有出人意料的联系（McCord 1966）。

\entry{9.4}{}设 $T$ 为拓扑空间，$\approx$ 为等价关系，$\pi:T\to T/{\approx}$ 为投影。规定 $A\subseteq T/{\approx}$ 为开集，当且仅当 $\pi^{-1}(A)$ 在 $T$ 中开，由此得到\textbf{商拓扑}。若 $S_i$（$i\in I$）是 $T$ 的两两不交子集，则 $T/(S_i)_{i\in I}$ 表示分别将每个 $S_i$ 中的点识别为一点所得的商空间。例如，
\[
\operatorname{cone}(T)=T\times[0,1]/(T\times\{1\})
\]
是 $T$ 上的\textbf{锥}，而
\[
\operatorname{susp}(T)=T\times[0,1]/(T\times\{0\},T\times\{1\})
\]
是 $T$ 的\textbf{悬挂}。$d$ 维闭球将其边界压成一点后同胚于 $d$ 维球面：$B^d/S^{d-1}\cong S^d$。

\sourcepage{1845}
若 $(T_i,x_i)_{i\in I}$ 为两两不交的带基点空间族，则其\textbf{楔和}为
\[
\bigvee_{i\in I}(T_i,x_i)=\left(\coprod_{i\in I}T_i\right)\Big/\left(\bigcup_{i\in I}\{x_i\}\right).
\]
两空间的\textbf{连接} $T_1*T_2$ 由 $T_1\times T_2\times[0,1]$ 取商得到：对每个 $t\in T_1$，将 $\{(t,x,0):x\in T_2\}$ 压成一点；对每个 $s\in T_2$，将 $\{(y,s,1):y\in T_1\}$ 压成一点。

两个顶点集不交的单纯复形的连接为
\[
\Delta_1*\Delta_2=\Delta_1\cup\Delta_2\cup\{\sigma\cup\tau:\sigma\in\Delta_1,\ \tau\in\Delta_2\}.
\]
若 $\Gamma_i$ 是有 $i$ 个顶点的零维复形（$i=1,2$），则 $\operatorname{cone}(\Delta)=\Delta*\Gamma_1$，$\operatorname{susp}(\Delta)=\Delta*\Gamma_2$。有同胚
\begin{equation}
\real{\Delta_1*\Delta_2}\cong\real{\Delta_1}*\real{\Delta_2}.\tag{9.5}
\end{equation}
若两个复形不局部有限，右侧的拓扑可能需要改成相应的紧生成拓扑，见 Walker (1988)。特别地，$\real{\operatorname{cone}(\Delta)}\cong\operatorname{cone}(\real\Delta)$，$\real{\operatorname{susp}(\Delta)}\cong\operatorname{susp}(\real\Delta)$。

连接还有如下几何实现。先在同一个足够高维的 $\R^d$ 中实现 $\Delta_1,\Delta_2$，使连接两侧点的任意两条不同线段，其内部不相交。再取所有这些线段的并，并赋予 $\R^d$ 的子空间拓扑，即得到 $\real{\Delta_1*\Delta_2}$。

复形 $\Delta$ 的\textbf{$p$ 重删除连接} $\Delta_*^{(p)}$ 定义如下：取不交副本 $\Delta_1,\ldots,\Delta_p$ 及同构 $f_i:\Delta_i\to\Delta$，从连接 $\Delta_1*\cdots*\Delta_p$ 中只保留满足
\[
f_i(\sigma_i)\cap f_j(\sigma_j)=\varnothing\quad(i\ne j)
\]
的面 $\sigma_1\cup\cdots\cup\sigma_p$。其组合应用见 Sarkaria (1990, 1991a,b) 及 Živaljević 与 Vrećica (1992)。

偏序集的\textbf{直积} $P\times Q$ 按逐坐标关系排序。它们的\textbf{连接}或\textbf{序和} $P*Q$ 是不交并，保留各自的序，并规定 $P$ 的每个元素都小于 $Q$ 的每个元素。显然 $\Delta(P*Q)=\Delta(P)*\Delta(Q)$。有以下同胚（Quillen 1978；Walker 1988）：
\begin{align}
\real{P\times Q}&\cong\real P\times\real Q,\tag{9.6}\\
\real{(P\times Q)_{>(x,y)}}&\cong\real{P_{>x}}*\real{Q_{>y}},\tag{9.7}\\
\real{((x,y),(x',y'))}&\cong\operatorname{susp}\big(\real{(x,x')}*\real{(y,y')}\big),\tag{9.8}
\end{align}
其中最后一式要求 $x<x'$ 于 $P$、$y<y'$ 于 $Q$。若相关序复形不局部有限，右侧拓扑仍须特别留意。

\entry{9.9}{}对 $\sigma\in\Delta\cup\{\varnothing\}$，定义以下子复形：
\begin{align*}
\operatorname{dl}_\Delta(\sigma)&=\{\tau\in\Delta:\tau\cap\sigma=\varnothing\} &&\text{（删除）},\\
\operatorname{st}_\Delta(\sigma)&=\{\tau\in\Delta:\tau\cup\sigma\in\Delta\} &&\text{（星形邻域）},\\
\lk_\Delta(\sigma)&=\{\tau\in\Delta:\tau\cap\sigma=\varnothing,\ \tau\cup\sigma\in\Delta\} &&\text{（链接）}.
\end{align*}
显然 $\operatorname{dl}(\sigma)\cap\operatorname{st}(\sigma)=\lk(\sigma)$，且 $\sigma*\lk(\sigma)=\operatorname{st}(\sigma)$。若 $\sigma$ 为顶点，另有 $\operatorname{dl}(\sigma)\cup\operatorname{st}(\sigma)=\Delta$。空面时，三者全为 $\Delta$。
\tnote{在 $\sigma*\lk(\sigma)$ 中，$\sigma$ 指由它的全部非空子集组成的单形。链接可能没有非空面，连接与空面有关的退化情形须采用相应的增广约定。}

\subsection*{同伦与同调}\addcontentsline{toc}{subsection}{同伦与同调}
\sourcepage{1846}
\entry{9.10}{}拓扑空间之间的两个映射 $f_0,f_1:T_1\to T_2$ 称为\textbf{同伦的}，记为 $f_0\sim f_1$，若存在映射 $F:T_1\times[0,1]\to T_2$，使对每个 $t\in T_1$，都有 $F(t,0)=f_0(t)$、$F(t,1)=f_1(t)$。$F$ 称为一个\textbf{同伦}。记住，拓扑空间之间的所有映射都假定连续。若存在 $f_1:T_1\to T_2$、$f_2:T_2\to T_1$，使
\[
f_2\circ f_1\sim\operatorname{id}_{T_1},\qquad
f_1\circ f_2\sim\operatorname{id}_{T_2},
\]
则称两个空间具有相同\textbf{同伦型}，或\textbf{同伦等价}，记为 $T_1\simeq T_2$。同伦等价于一点的空间称为\textbf{可缩的}。

标准 $(d-1)$ 维球面与 $d$ 维闭球分别为
\[
S^{d-1}=\{x\in\R^d:\|x\|=1\},\qquad
B^d=\{x\in\R^d:\|x\|\le1\}.
\]
注意 $S^{-1}=\varnothing$，$S^0$ 为两点，$B^0$ 为一点。球面与闭球对连接运算封闭（在同胚意义下）：
\[
S^a*S^b\cong S^{a+b+1},\qquad
B^a*B^b\cong B^a*S^b\cong B^{a+b+1}.
\]

空间 $T$ 称为\textbf{$k$ 连通的}，若对全部 $0\le i\le k$，每个映射 $f:S^i\to T$ 都可延拓为 $\widehat f:B^{i+1}\to T$，即 $\widehat f$ 在 $S^i$ 上等于 $f$。特别地，0 连通就是道路连通。$k$ 连通性是同伦不变量，即会传给同伦等价的空间。$S^d$ 是 $(d-1)$ 连通的，但不是 $d$ 连通的（见定理 13.1），而 $B^d$ 可缩。方便起见约定：$(-1)$ 连通就是非空；$k\le-2$ 时任何空间，无论是否为空，都为 $k$ 连通。

单纯复形 $\Delta$ 可缩，当且仅当它对所有 $k\ge0$ 都 $k$ 连通；等价地，只须对 $0\le k\le\dim\Delta$ 成立。一般拓扑空间的相应命题在非平凡的那个方向并不成立。此外，复形 $k$ 连通当且仅当它的 $(k+1)$ 骨架 $k$ 连通。

设 $\pi_i(T)=\pi_i(T,x)$ 为从带基点的 $i$ 维球面到带基点空间 $(T,x)$ 的映射的同伦类集，其中 $i\ge0$，且 $f((1,0,\ldots,0))=x$。当 $i\ge1$ 时，存在使它成为群的合成运算，所得群称为 $T$ 在 $x$ 处的\textbf{第 $i$ 同伦群}；$i\ge2$ 时该群交换。$\pi_1(T)$ 是\textbf{基本群}；$\pi_1(T)=0$ 时称空间\textbf{单连通}。空间 $T$ 为 $k$ 连通，当且仅当对所有 $0\le i\le k$ 以及 $x\in T$，都有 $\pi_i(T,x)=0$。因此 1 连通意味着单连通且道路连通。
\tnote{原文把“基本群平凡”与“道路连通”分开陈述。许多教材把道路连通也包含在“单连通”的定义内，阅读时应留意这一约定。}

\entry{9.11}{}复形 $\Delta$ 以交换群 $G$ 为系数的\textbf{单纯同调群} $H_i(\Delta,G)$ 和\textbf{约化单纯同调群} $\widetilde H_i(\Delta,G)$ 的定义，参阅 Munkres (1984a) 或 Spanier (1966)。记 $\widetilde H_i(\Delta)=\widetilde H_i(\Delta,\Z)$。须注意退化情形：
\[
\widetilde H_i(\varnothing)\cong
\begin{cases}\Z,&i=-1,\\0,&i\ne-1.\end{cases}
\]
若 $\Delta\ne\varnothing$，则 $i<0$ 或 $i>\dim\Delta$ 时 $\widetilde H_i(\Delta)=0$，而 $\widetilde H_0(\Delta)\cong\Z^{c-1}$，其中 $c$ 是连通分支数。对 $i\ne-1,0$，$H_i(\Delta)=\widetilde H_i(\Delta)$；另有 $H_{-1}(\Delta)=0$ 和 $H_0(\Delta)\cong\widetilde H_0(\Delta)\oplus\Z$。

\sourcepage{1847}
设 $\Delta_1,\Delta_2$ 为有限复形，并假定 $p+q=i-1$ 时，$\widetilde H_p(\Delta_1)$ 与 $\widetilde H_q(\Delta_2)$ 至少一个无挠。则
\begin{equation}
\widetilde H_{i+1}(\Delta_1*\Delta_2)
\cong\bigoplus_{p+q=i}\big(\widetilde H_p(\Delta_1)\otimes\widetilde H_q(\Delta_2)\big).
\tag{9.12}
\end{equation}
以域为系数的约化同调无需任何限制，也有同样分解。详见 Milnor (1956)，或 Cooke 与 Finney (1967) 第 V 章。

对有限单纯复形 $\Delta$，令 $\beta_i:=\operatorname{rank}H_i(\Delta)=\dim_\Q H_i(\Delta,\Q)$，$i\ge0$。\textbf{Betti 数} $\beta_i$ 满足 \textbf{Euler--Poincaré 公式}：
\begin{equation}
\sum_{i\ge0}(-1)^i|\Delta^i|=\sum_{i\ge0}(-1)^i\beta_i.\tag{9.13}
\end{equation}
两侧任一边都可用来定义\textbf{Euler 示性数} $\chi(\Delta)$。\textbf{约化 Euler 示性数}为 $\widetilde\chi(\Delta)=\chi(\Delta)-1$。对任意域 $\mathbf k$，取 $\beta_i=\dim_{\mathbf k}H_i(\Delta,\mathbf k)$，公式仍成立，尽管各个 $\beta_i$ 本身可能依赖系数域。面数 $f_i=|\Delta^i|$ 与 Betti 数间还有其他关系（Björner 与 Kalai 1988）。对于各种特殊复形类，关于 $f$ 向量 $f(\Delta)=(f_0,f_1,\ldots)$ 已知很多结果。多胞形边界这一重要情形见 Klee 与 Kleinschmidt 所写的第 18 章，更一般复形类的综述见 Björner 与 Kalai (1989)。

局部有限偏序集的 Möbius 函数在 Gessel 与 Stanley 所写的第 21 章定义。利用 (9.13)，该章源于 P. Hall 的定理 13.4 可写成
\begin{equation}
\mu(x,y)=\widetilde\chi\big(\Delta((x,y))\big),\qquad x<y.\tag{9.14}
\end{equation}
右侧是开区间 $(x,y)$ 的序复形的约化 Euler 示性数。Möbius 函数与拓扑之间的这一联系，最早由 Rota (1964) 和 Folkman (1966) 指出，有许多有趣的延伸。

\entry{9.15}{}同伦等价的两个复形在各维具有同构的同调群。若对所有 $i\le k$，都有 $\widetilde H_i(\Delta,G)=0$，称 $\Delta$ \textbf{在 $G$ 上 $k$ 无环}。因此 $(-1)$ 无环表示非空，0 无环表示非空且连通。若对全部 $i\in\Z$ 都有 $\widetilde H_i(\Delta,G)=0$，则称 $\Delta$ \textbf{在 $G$ 上无环}，或称\textbf{$G$ 无环}。省略系数群时，一律指 $G=\Z$。
\tnote{“无环”（acyclic）在这里指约化同调消失；不是说复形的一维骨架没有图论意义的圈。}

下面列出复形的同伦性质与同调之间的一些常用联系，它们来自 Hurewicz 与 Whitehead 定理，见 Spanier (1966)。

\entry{9.16}{}对 $k\ge1$，$\Delta$ 为 $k$ 连通，当且仅当它在 $\Z$ 上 $k$ 无环且单连通。
\sourcepage{1848}
\entry{9.17}{}$\Delta$ 可缩，当且仅当它 $\Z$ 无环且单连通。

\entry{9.18}{}若 $\Delta$ 单连通，$d>1$，对所有 $i\ne d$ 有 $\widetilde H_i(\Delta)=0$，且 $\widetilde H_d(\Delta)\cong\Z^k$，则 $\Delta$ 同伦等价于 $k$ 个 $d$ 维球面的楔和。

\entry{9.19}{}设 $\dim\Delta=d\ge0$。则 $\Delta$ 为 $(d-1)$ 连通，当且仅当它同伦等价于若干 $d$ 维球面的楔和。

\textbf{原文附注：}对不能三角剖分的空间，9.17--9.19 的类似命题可能不成立。

\entry{9.20}{}若 $\Delta_1$ 为 $k_1$ 无环，$\Delta_2$ 为 $k_2$ 无环，则 $\Delta_1*\Delta_2$ 为 $(k_1+k_2+2)$ 无环。这来自 (9.12)。再利用 9.16 可得：若 $\Delta_i$ 为 $k_i$ 连通，则 $\Delta_1*\Delta_2$ 为 $(k_1+k_2+2)$ 连通。另见 Milnor (1956)。


\omitsection{10}{组合同伦定理（Combinatorial homotopy theorems）}
\setcounter{section}{10}
\section{具有特殊结构的复形}
\sourcepage{1853}
本节回顾在组合学中经常遇到的一些复形性质，它们表达了结构在某种意义上的简单性。
\subsection*{可塌缩与可壳复形}\addcontentsline{toc}{subsection}{可塌缩与可壳复形}
\entry{11.1}{}设 $\Delta$ 为单纯复形，且 $\sigma\in\Delta$ 恰为一个单形 $\tau\in\Delta$ 的真面。则
\[
\Delta'=\Delta\setminus\{\sigma,\tau\}
\]
称为由 $\Delta$ 经一次\textbf{初等塌缩}得到，而 $\Delta$ 由 $\Delta'$ 经一次\textbf{初等反塌缩}得到。注意 $\Delta'\simeq\Delta$。若经过一列初等塌缩，能将 $\Delta$ 化为单点，则称它\textbf{可塌缩}。

\textbf{非回避复形}（nonevasive complex）递归定义如下：单个顶点是非回避的；若某个顶点 $x\in\Delta^0$ 的链接 $\lk_\Delta(x)$ 与删除 $\operatorname{dl}_\Delta(x)$ 都非回避，则 $\Delta$ 非回避。以下蕴含都是严格的，即逆命题不成立：
\[
\text{锥}\Longrightarrow\text{非回避}\Longrightarrow\text{可塌缩}
\Longrightarrow\text{可缩}\Longrightarrow\Z\text{ 无环}.
\]
此外，对任意域 $\mathbf k$：
\[
\Z\text{ 无环}\Longrightarrow\mathbf k\text{ 无环}
\Longrightarrow\Q\text{ 无环}\Longrightarrow\widetilde\chi=0,
\]
并且，$\Z$ 无环当且仅当对每个素数 $p$ 都 $\Z_p$ 无环。

Kahn 等人 (1984) 定义非回避复形，是为了刻画第 2 节讨论的自变量复杂度。复形 $\Delta$ 非回避，当且仅当对所有 $F\subseteq\Delta^0$，都能用少于 $|\Delta^0|$ 个“$x\in F$ 吗？”这样的询问，判定 $F\in\Delta$ 是否成立。
\tnote{这个询问模型按组合学约定把空面也纳入，见 9.1 的原文附注；否则单个顶点的基例会与空集输入产生约定冲突。}

可塌缩性在组合拓扑中已有长期研究。值得注意：两个单连通有限复形同伦等价，当且仅当能通过一列初等塌缩与反塌缩将一个变为另一个，见 Cohen (1973)。特别地，可缩复形恰为能通过塌缩和反塌缩化为一点的复形。

\sourcepage{1854}
若偏序集 $P$ 的元素 $x$ 满足 $P_{>x}$ 有最小元或 $P_{<x}$ 有最大元，称 $x$ \textbf{不可约}。若有限偏序集可以反复移去不可约元素，最后成为单元素偏序集，称其\textbf{可拆除}。可拆除偏序集是非回避的。Stong (1966) 给出的可拆除偏序集的拓扑刻画已在 9.3 提到。\textbf{有向偏序集}，即任意 $x,y\in P$ 都有共同上界 $z\in P$，是可缩的。

\entry{11.2}{}设 $\Delta$ 为纯 $d$ 维复形，且 $k$ 维面 $\sigma$ 恰包含于一个 $d$ 维面 $\tau$。则
\[
\Delta'=\Delta\setminus\{\gamma:\sigma\subseteq\gamma\subseteq\tau\}
\]
称为由 $\Delta$ 经一次\textbf{$(k,d)$ 塌缩}得到。若 $\sigma\ne\tau$，则 $\Delta'\simeq\Delta$。若可以通过一列 $(k,d)$ 塌缩（$0\le k\le d$）将 $\Delta$ 化为一个 $d$ 单形，则称 $\Delta$ \textbf{可壳}。
\tnote{这里允许 $k=d$，即直接删除一个极大面；这种步骤未必保持同伦型。因此“可壳”不蕴含“可塌缩”或“可缩”。}

纯单纯复形称为\textbf{顶点可分解}，若满足下列之一：(i) $\Delta=\varnothing$；(ii) $\Delta$ 为单个顶点；(iii) 存在 $x\in\Delta^0$，使 $\lk_\Delta(x)$ 与 $\operatorname{dl}_\Delta(x)$ 都顶点可分解。例如，每个单形及其边界都顶点可分解。\textbf{可构造复形}递归定义为：(i) 每个单形及空复形均可构造；(ii) 若 $\Delta_1,\Delta_2,\Delta_1\cap\Delta_2$ 均可构造，且
\[
\dim\Delta_1=\dim\Delta_2=1+\dim(\Delta_1\cap\Delta_2),
\]
则 $\Delta_1\cup\Delta_2$ 可构造。纯 $d$ 维复形的下列蕴含是严格的：
\[
\begin{gathered}
\text{顶点可分解}\Longrightarrow\text{可壳}\Longrightarrow\text{可构造}\\
\Longrightarrow(d-1)\text{ 连通}.
\end{gathered}
\]
第一条蕴含以及顶点可分解的定义来自 Provan 与 Billera (1980)。可壳性的历史可追溯到 19 世纪，见 Grünbaum (1967)。可构造复形由 M. Hochster 定义，见 Stanley (1977)。

可壳性通常被看作把复形\emph{拼起来}的方式，因此以下等价定义更常用。有限纯 $d$ 维复形可壳，当其 $d$ 维面可排为 $\sigma_1,\ldots,\sigma_t$，使对 $2\le k\le t$，
\[
(\partial\sigma_1\cup\cdots\cup\partial\sigma_{k-1})\cap\partial\sigma_k
\]
为纯 $(d-1)$ 维复形；这里 $\partial\sigma_j=2^{\sigma_j}\setminus\{\varnothing,\sigma_j\}$ 是该单形的边界复形。等价地，对任意 $1\le i<k\le t$，存在 $j<k$ 使
\[
\sigma_i\cap\sigma_k\subseteq\sigma_j\cap\sigma_k,\qquad
\dim(\sigma_j\cap\sigma_k)=d-1.
\]
用语言说，第 $k$ 个极大面与先前极大面的并，沿着其边界的一部分相交，而这部分是该单形若干极大真面的并。这样的极大面排序称为\textbf{壳序}。

若 $\Delta$ 可壳或可构造，则每个链接 $\lk_\Delta(\sigma)$ 也有相同性质。另一些复形与偏序集构造也保持可壳性，例如定理 11.13。单纯可壳性的一些基本性质，包括无限复形情形，综述于 Björner (1984b)。胞腔复形的可壳性见 Danaraj 与 Klee (1974)、Björner (1984a)，以及 Klee 与 Kleinschmidt 所写的第 18 章。证明偏序集的序复形可壳，有一种专门方法称为\textbf{字典序可壳性}，详见 Björner (1980) 和 Björner 与 Wachs (1983, 1994)。可壳性、顶点可分解性及其大多数有用性质，都容易推广到非纯复形，见 Björner 与 Wachs (1994)。

\sourcepage{1855}
\entry{11.3}{}单纯 PL 球面与 PL 闭球的定义见 12.2，PL 表示分片线性（piecewise linear）。PL 性是一种组合性质：几何单纯复形是否为 PL，只取决于其抽象单纯复形。

为了证明具体复形同胚于球面或闭球，下列结果常常很有用。
\entry{定理 11.4}{}设 $\Delta$ 为可构造的 $d$ 维单纯复形。
\begin{enumerate}[label=(\roman*)]
\item 若每个 $(d-1)$ 维面恰包含于两个 $d$ 维面，则 $\Delta$ 是 PL 球面。
\item 若每个 $(d-1)$ 维面包含于一个或两个 $d$ 维面，且至少有一个只包含于一个 $d$ 维面，则 $\Delta$ 是 PL 闭球。
\end{enumerate}
这来自 PL 拓扑的一些基本事实，例如 12.2 引用的结果。对于可壳复形，这个定理隐含于 Bing (1964)，明确出现于 Danaraj 与 Klee (1974)。

若 $\Delta$ 是 $d$ 维球面（或任意无边界流形）的三角剖分，$\sigma\in\Delta^k$，则 $\lk_\Delta(\sigma)$ 的同调与 $S^{d-1-k}$ 相同。若 $\sigma$ 为顶点，链接甚至与 $S^{d-1}$ 同伦等价。然而，若 $\Delta$ 是 PL $d$ 维球面，则每个 $k$ 维面的链接本身就是 PL $(d-1-k)$ 维球面。
\tnote{原文这一段写“任意流形”；球面型局部同调的说法适用于无边界流形的面，或有边界流形中不位于边界的面。边界面需作相应调整。}

\subsection*{Cohen--Macaulay 复形}\addcontentsline{toc}{subsection}{Cohen--Macaulay 复形}
\entry{11.5}{}设 $\mathbf k$ 为域或整数环 $\Z$。有限维单纯复形 $\Delta$ 称为\textbf{在 $\mathbf k$ 上 Cohen--Macaulay}，记为 $\cm/\mathbf k$；当系数已明确或不重要时简称 CM，若对每个 $\sigma\in\Delta\cup\{\varnothing\}$，链接 $\lk_\Delta(\sigma)$ 在 $\mathbf k$ 上均为
\[
(\dim\lk_\Delta(\sigma)-1)\text{ 无环}.
\]
若对所有这样的 $\sigma$，链接均为 $(\dim\lk_\Delta(\sigma)-1)$ 连通，则称 $\Delta$ 为\textbf{同伦 Cohen--Macaulay}（homotopy-CM）。以下蕴含为严格蕴含：
\[
\begin{gathered}
\text{可构造}\Longrightarrow\text{同伦 CM}\Longrightarrow\cm/\Z\\
\Longrightarrow\cm/\mathbf k\Longrightarrow\cm/\Q,
\end{gathered}
\]
其中 $\mathbf k$ 为任意域。此外，$\cm/\Z$ 当且仅当对每个素数 $p$ 都为 $\cm/\Z_p$。第一条蕴含来自可构造性蕴含 $(d-1)$ 连通，且会被链接继承；第二条来自 9.15；其余来自万有系数定理。特别地，可壳复形都是同伦 CM 的。
\tnote{链中的严格性应按不同系数情形理解：特征为 0 的域与 $\Q$ 在这里给出相同的同调消失条件，因此不能把最后一条对每个固定域都理解成真包含。正文保留原文的概括表述。}

有限 CM 复形的一个重要方面，是存在等价的环论定义。设 $\Delta^0=\{x_1,\ldots,x_n\}$，在多项式环 $\mathbf k[x_1,\ldots,x_n]$ 中，取由所有满足
\[
\{x_{i_1},\ldots,x_{i_s}\}\notin\Delta,\qquad
1\le i_1<\cdots<i_s\le n,\quad s\ge1
\]
的单项式 $x_{i_1}\cdots x_{i_s}$ 生成的理想 $I$。令
\[
\mathbf k[\Delta]=\mathbf k[x_1,\ldots,x_n]/I,
\]
称为 $\Delta$ 的 \textbf{Stanley--Reisner 环}或\textbf{面环}。则 $\Delta$ 为 $\cm/\mathbf k$，当且仅当 $\mathbf k[\Delta]$ 在交换代数意义下为 Cohen--Macaulay 环（Reisner 1976）。

\sourcepage{1856}
单纯复形的环论方面及其组合应用，可参阅 Stanley (1983a)。该书还讨论了其他由环论引出的复形类，例如 Gorenstein 复形和 Buchsbaum 复形。复形环论及 Reisner 定理的其他处理见 Baclawski 与 Garsia (1981)、Yuzvinsky (1987)。另见第 41 章“纯数学中的组合学”的第 5 节。

Cohen--Macaulay 复形与偏序集约于 1974--1975 年在 Baclawski (1976, 1980)、Hochster (1977)、Reisner (1976)、Stanley (1975, 1977) 的工作中引入。同伦 CM 的概念最早出现于 Quillen (1978)。Björner、Garsia 与 Stanley (1982) 给出 CM 偏序集的初等介绍。

Cohen--Macaulay 性的一个突出组合应用，是 Stanley 证明了 $d$ 维球面具有 $n$ 个顶点的三角剖分中，各维面数的精确上界（Stanley 1975, 1983a；另见 Klee 与 Kleinschmidt 所写的第 18 章）。关于下界的一个应用见 Stanley (1987a)。

\entry{11.6}{}纯 $d$ 维复形称为\textbf{强连通}或\textbf{对偶连通}，若每对极大面 $\sigma,\tau\in\Delta^d$ 都可由一列极大面
\[
\sigma=\sigma_0,\sigma_1,\ldots,\sigma_n=\tau
\]
连接，使对每个 $1\le i\le n$，有 $\dim(\sigma_{i-1}\cap\sigma_i)=d-1$。

\entry{命题 11.7}{}每个 CM 复形都是纯的且强连通。

这来自以下引理；对 $\dim\Delta$ 归纳即可证明：设 $\Delta$ 为有限维单纯复形，假定对每个满足 $\dim\lk_\Delta(\sigma)\ge1$ 的 $\sigma\in\Delta\cup\{\varnothing\}$，$\lk_\Delta(\sigma)$ 都连通。则 $\Delta$ 纯且强连通。

CM 性是拓扑不变量：$\Delta$ 是否为 $\cm/\mathbf k$ 只取决于 $\real\Delta$ 的拓扑。这由 Munkres (1984b) 给出的下列等价表述推出。

\entry{定理 11.8}{}有限维复形 $\Delta$ 为 $\cm/\mathbf k$，当且仅当其空间 $T=\real\Delta$ 满足：对所有 $p\in T$ 及 $i<\dim\Delta$，
\[
\widetilde H_i(T)=H_i(T,T\setminus\{p\})=0.
\]
此处 $\widetilde H_i$ 为约化奇异同调，$H_i$ 为相对奇异同调，系数均取 $\mathbf k$。其推论是：若可三角剖分流形 $M$（可有边界）满足 $\widetilde H_i(M)=0$ 对全部 $i<\dim M$ 成立，则 $M$ 的每个三角剖分都 CM。例如：(1) $d$ 维球面、闭球或 $\R^d$ 的每个三角剖分均为 $\cm/\Z$，但未必同伦 CM；注意，同伦 CM 不是拓扑不变量。(2) 实射影 $d$ 维空间的三角剖分为 $\cm/\mathbf k$，当且仅当 $\operatorname{char}\mathbf k\ne2$。
\tnote{第二个例子在 $d\ge2$ 的通常语境下成立；$\R P^1\cong S^1$ 对任意系数域均 CM，故低维情形须单独看待。}

\entry{11.9}{}Cohen--Macaulay 偏序集，即序复形为 CM 的偏序集，值得补充一点说明。设 $P$ 具有有限秩，$\sigma:x_0<\cdots<x_k$ 为其中一条链。它在序复形中的链接是以下序和的序复形：
\[
P_{<x_0}*(x_0,x_1)*\cdots*(x_{k-1},x_k)*P_{>x_k}.
\]
因此，由 9.20，$P$ 为 CM（相应地，同伦 CM），当且仅当 $\widehat P$ 的每个开区间 $(x,y)$ 都为 $(\rk(x,y)-1)$ 无环（相应地，$(\rk(x,y)-1)$ 连通）。

\sourcepage{1857}
Cohen--Macaulay 偏序集在交换代数中的一些应用，见第 41 章“纯数学中的组合学”的第 5 节。

\entry{11.10}{}组合学中有大量可壳复形与 CM 复形。这里只列举少数重要例子。
\begin{enumerate}[label=(\roman*)]
\item 单纯凸多胞形的边界复形可壳（Bruggesser 与 Mani 1971；Danaraj 与 Klee 1974；另见 Klee 与 Kleinschmidt 所写的第 18 章）。每个单纯 PL 球面都是某个可壳闭球的边界（Pachner 1986）。存在三维闭球的不可壳三角剖分（M. E. Rudin），也存在三维球面的不可壳三角剖分（见下）。球面与闭球可壳性的综述见 Danaraj 与 Klee (1978)。
\item 对任意单纯球面，均有
\[
\text{可构造}\Longrightarrow\mathrm{PL}\Longrightarrow\text{同伦 CM}.
\]
五维球面有非同伦 CM 的三角剖分（R. D. Edwards，见 Daverman 1986），也有不可构造的 PL 三角剖分（Mandel 1982）。三维球面的每个三角剖分都是 PL，但并非全部可壳（Lickorish 1991；另见 Vince 1985）。正则复多胞形的面格是 CM 的（Orlik 1990）。
\item 拟阵的独立集复形可构造（Stanley 1977），且顶点可分解（Provan 与 Billera 1980）。更一般地，由贪心结构（greedoid）的基的补集所生成的复形顶点可分解（Björner、Korte 与 Lovász 1985）。拟阵产生的复形见 Björner (1992)。
\item 每个有限秩半模格，特别是几何格或模格，都是 CM 的（Folkman 1966）且可壳（Björner 1980）。对几何格 $L$ 的任意元素 $x$，偏序集 $L\setminus[x,\hat1]$ 可壳（Wachs 与 Walker 1986）。
\item Tits 建筑是 CM 的（Solomon--Tits，见 Brown 1989 或 Ronan 1989）且可壳（Björner 1984b）。Ronan (1981)、Smith (1988)、Tits (1981) 等人，为了在群论中应用，也研究了与群有关的更一般几何的拓扑。通论见 Buekenhout (1995)、Ronan (1989)。
\item Quillen (1978) 证明，有限群的初等交换 $p$ 子群偏序集在某些情形下为同伦 CM；另见 Stong (1984)。有限群 $G$ 的整个子群格可壳（或 CM），当且仅当 $G$ 超可解（Björner 1980）。各种子群偏序集已从拓扑角度得到研究，相关文献指南见 Thévenaz (1987)、Webb (1987)、Welker (1994)。
\end{enumerate}

\subsection*{诱导子复形}\addcontentsline{toc}{subsection}{诱导子复形}
连通性、Cohen--Macaulay 性等，在一定条件下会被适当子复形继承。对于单纯复形 $\Delta$ 及 $A\subseteq\Delta^0$，记
\[
\Delta_A=\{\sigma\in\Delta:\sigma\subseteq A\},
\]
即 $A$ 上的\textbf{诱导子复形}。

\entry{引理 11.11}{}设 $\Delta$ 为有限维复形，$A\subseteq V=\Delta^0$。假定对每个 $\sigma\in\Delta_{V\setminus A}$，$\lk_\Delta(\sigma)$ 都为 $k$ 连通。则 $\Delta_A$ 为 $k$ 连通，当且仅当 $\Delta$ 为 $k$ 连通。

\sourcepage{1858}
\entry{引理 11.12}{}设 $P$ 为有限秩偏序集，$A\subseteq P$。假定对每个 $x\in P\setminus A$，$P_{>x}$ 都为 $k$ 连通。则 $A$ 为 $k$ 连通，当且仅当 $P$ 为 $k$ 连通。

\begin{proofcn}
这两个引理等价。先证明引理 11.12。令 $f:A\to P$ 为嵌入映射。对 $x\in P$，
\[
f^{-1}(P_{\ge x})=
\begin{cases}
A_{\ge x},&x\in A,\\
P_{>x}\cap A,&x\notin A.
\end{cases}
\]
其中 $A_{\ge x}$ 是锥，故可缩；而 $P_{>x}\cap A$ 由对 $\rk P$ 的归纳为 $k$ 连通。于是由定理 10.5(ii) 得结论。

为证明引理 11.11，令 $P=\mathcal P(\Delta)$，
\[
Q=\{\tau\in\Delta:\tau\cap A\ne\varnothing\}\subseteq P.
\]
对 $\sigma\in P\setminus Q$，有 $P_{>\sigma}\cong\mathcal P(\lk_\Delta(\sigma))$，故为 $k$ 连通，可以应用引理 11.12。另一方面，由推论 10.12，$Q$ 上的映射 $f(\tau)=\tau\cap A$ 诱导 $Q$ 与 $f(Q)=\mathcal P(\Delta_A)$ 之间的同伦等价。
\end{proofcn}

始终用 $k$ 无环替换 $k$ 连通，即得到两个引理的同调版本，证明方法平行。若将假设中的“$k$ 连通”改为“可缩”，则结论为 $\Delta_A$ 与 $\Delta$ 同伦等价，相应地 $A$ 与 $P$ 同伦等价。

\entry{定理 11.13}{}设 $\Delta$ 为纯 $d$ 维单纯复形，$A\subseteq\Delta^0$，$1\le m\le d$。假定每个极大面 $\sigma\in\Delta^d$ 都满足 $|A\cap\sigma|=m$。若 $\Delta$ 为 $\cm/\mathbf k$、同伦 CM 或可壳，则 $\Delta_A$ 继承同一性质。

关于 CM 性，这个结果的不同一般性版本分别由 Baclawski (1980)、Munkres (1984b)、Stanley (1979)、Walker (1981a) 证明。它容易由引理 11.11 得到。可壳性的证明见 Björner (1980, 1984b)。

设纯 $d$ 维复形 $\Delta$ 有映射 $t:\Delta^0\to\{0,1,\ldots,d\}$，其在每个极大面上的限制均为双射。则称 $\Delta$ 为\textbf{完全平衡的}，或称有编号、有着色的复形，$t$ 称为\textbf{类型映射}。例如，纯偏序集的序复形完全平衡，以秩为类型映射，参见 9.2；建筑式关联几何（Buekenhout 1995）也产生完全平衡复形。这类 CM 复形由 Stanley (1979) 等人研究。

对每个 $J\subseteq\{0,1,\ldots,d\}$，\textbf{类型选择子复形}
\[
\Delta_{(J)}:=\Delta_{t^{-1}(J)}
\]
是在 $t^{-1}(J)$ 上的诱导子复形。定理 11.13 表明，若 $\Delta$ 为 CM，则 $\Delta_{(J)}$ 也 CM，从而为 $(|J|-2)$ 无环。某种逆命题也成立，即下列结果；它给出完全平衡复形的 CM 性的另一种刻画。有限 CM 情形由 Baclawski 与 Garsia (1981) 给出；包括同伦情形的一般版本由 J. Walker 于 1981 年致作者的信中给出。

\entry{定理 11.14}{}设 $\Delta$ 为纯 $d$ 维完全平衡复形。则 $\Delta$ 为 $\cm/\mathbf k$（相应地，同伦 CM），当且仅当对每个 $J\subseteq\{0,1,\ldots,d\}$，$\Delta_{(J)}$ 在 $\mathbf k$ 上为 $(|J|-2)$ 无环（相应地，为 $(|J|-2)$ 连通）。


\omitsection{12}{胞腔复形（Cell complexes）}
\noindent 原文未编号分题：多面体复形与 PL 拓扑；正则胞腔复形。\par
\omitsection{13}{不动点与对径性定理（Fixed-point and antipodality theorems）}
\blankrefs
\end{document}
